\documentclass[main.tex]{subfiles}
\newcommand{\DiagramaEscala}{0.9}
\begin{document}

\section{Ecuaciones Diferenciales Ordinarias}

%--------------------------------------------------
% Diapositiva 1
%--------------------------------------------------
%--- Diapositiva: Definici\'on de EDO ---
\begin{frame}{\small Definici\'on de Ecuaci\'on Diferencial Ordinaria (EDO)}
\begin{itemize}
  \item Una \textbf{Ecuaci\'on Diferencial Ordinaria (EDO)} contiene una funci\'on desconocida y sus derivadas respecto a una sola variable independiente.
  \item Ejemplo (EDO de primer orden):
  \[
    \frac{dv}{dt} = \frac{g c}{m} - \frac{c}{m}v
  \]
  \item La variable \textbf{dependiente} es $v(t)$ (velocidad), y la \textbf{independiente} es $t$ (tiempo).
  \item Este tipo de ecuaci\'on se conoce tambi\'en como ecuaci\'on de \emph{raz\'on de cambio}.
\end{itemize}
\end{frame}

%--- Diapositiva: Clasificaci\'on por orden ---
\begin{frame}{\small Clasificaci\'on seg\'un el orden de la EDO}
\begin{itemize}
  \item El \textbf{orden} de una EDO est\'a determinado por la derivada de mayor orden presente.
  \item Ejemplo de \textbf{primer orden}: \( \frac{dv}{dt} = \cdots \)
  \item Ejemplo de \textbf{segundo orden}:
  \[
    m \frac{d^2x}{dt^2} + c \frac{dx}{dt} + kx = 0
  \]
  \item Este tipo de ecuaci\'on representa un sistema masa-resorte con amortiguamiento.
\end{itemize}
\end{frame}

%--- Diapositiva: Reducci\'on a sistema de primer orden ---
\begin{frame}{\small Reducci\'on de EDO de segundo orden a sistema de primer orden}
\begin{itemize}
  \item Se define una nueva variable auxiliar:
  \[
    y = \frac{dx}{dt}
  \]
  \item Derivando:
  \[
    \frac{dy}{dt} = \frac{d^2x}{dt^2}
  \]
  \item Sustituyendo en la ecuaci\'on original:
  \[
    m \frac{dy}{dt} + c y + k x = 0
  \]
\end{itemize}
\end{frame}

%--- Diapositiva: Sistema equivalente ---
\begin{frame}{\small Sistema de ecuaciones equivalentes de primer orden}
\begin{align*}
  \frac{dx}{dt} &= y \\
  \frac{dy}{dt} &= -\frac{c}{m} y - \frac{k}{m} x
\end{align*}
\begin{itemize}
  \item Estas dos ecuaciones forman un sistema de primer orden equivalente a la EDO de segundo orden original.
  \item Este m\'etodo es muy utilizado para resolver num\'ericamente EDOs de orden superior.
\end{itemize}
\end{frame}


%--- Diapositiva: Resoluci\'on anal\'itica sin computadora ---
\begin{frame}{\small Resoluci\'on de EDO sin computadora}
\begin{itemize}
  \item Algunas EDO pueden resolverse por \textbf{integraci\'on anal\'itica}:
  \[
    dv = \left(\frac{gc}{m} - \frac{c}{m}v\right)dt
  \]
  \item Integrando ambos lados (integral indefinida):
  \[
    \int dv = \int \left(\frac{gc}{m} - \frac{c}{m}v\right) dt
  \]
  \item Esta t\'ecnica solo es posible si la integral puede resolverse exactamente.
\end{itemize}
\end{frame}

%--- Diapositiva: Soluci\'on del paracaidista ---
\begin{frame}{\small Soluci\'on anal\'itica del paracaidista}
\begin{itemize}
  \item Para el problema del paracaidista (con $v=0$ en $t=0$):
  \[
    v(t) = \frac{gm}{c}(1 - e^{-\frac{c}{m}t})
  \]
  \item Esta es una \textbf{soluci\'on exacta} que se obtiene al resolver la integral anal\'itica.
  \item Sin embargo, no todas las EDO tienen soluciones anal\'iticas.
\end{itemize}
\end{frame}

%--- Diapositiva: Linealizaci\'on ---
\begin{frame}{\small Linealizaci\'on de EDO no lineales}
\begin{itemize}
  \item Las EDO lineales tienen la forma:
  \[
    a_n(x) y^{(n)} + \dots + a_1(x) y' + a_0(x) y = f(x)
  \]
  \item Son m\'as f\'aciles de resolver anal\'iticamente.
  \item Las EDO no lineales frecuentemente se \textbf{linealizan} para su estudio.
\end{itemize}
\end{frame}

%--- Diapositiva: Ejemplo del p\'endulo ---
\begin{frame}{\small Ejemplo: Movimiento de un p\'endulo}
\begin{itemize}
  \item EDO original (no lineal):
  \[
    \frac{d^2\theta}{dt^2} + \frac{g}{l} \sin\theta = 0
  \]
  \item Para peque\~nos \( \theta \):
  \[
    \sin\theta \approx \theta \quad \Rightarrow \quad \frac{d^2\theta}{dt^2} + \frac{g}{l} \theta = 0
  \]
  \item Resultado: EDO lineal que puede resolverse anal\'iticamente.
\end{itemize}
\end{frame}

%--- Diapositiva: Limitaciones de la linealizaci\'on ---
\begin{frame}{\small Limitaciones de la linealizaci\'on}
\begin{itemize}
  \item Linealizaci\'on solo es v\'alida para \textbf{peque\~nos desplazamientos} o simplificaciones.
  \item Para situaciones m\'as realistas (grandes desplazamientos):
  \begin{itemize}
    \item No es aplicable.
    \item Se requieren \textbf{m\'etodos num\'ericos}.
  \end{itemize}
  \item Las computadoras modernas permiten resolver EDOs complejas sin recurrir a simplificaciones.
\end{itemize}
\end{frame}

%--- Diapositiva: Introducci\'on a los m\'etodos de Runge-Kutta ---
\begin{frame}{\small M\'etodos de Runge-Kutta}
\begin{itemize}
  \item Consideramos EDO de la forma:
  \[
    \frac{dy}{dx} = f(x, y)
  \]
  \item Todos los m\'etodos de un paso siguen la estructura:
  \[
    y_{i+1} = y_i + \phi h
  \]
  \item Lo que var\'ia es \textbf{c\'omo se estima la pendiente} $\phi$.
  \item El m\'etodo m\'as simple es el \textbf{m\'etodo de Euler}.
\end{itemize}
\end{frame}

%--- Diapositiva: M\'etodo de Euler ---
\begin{frame}{\small M\'etodo de Euler}
\begin{itemize}
  \item Se estima la pendiente en $x_i$:
  \[
    \phi = f(x_i, y_i)
  \]
  \item Se aplica:
  \[
    y_{i+1} = y_i + f(x_i, y_i) h
  \]
  \item Es un m\'etodo simple pero con errores relativamente altos si $h$ no es peque\~no.
\end{itemize}
\end{frame}

%--- Diapositiva: Ejemplo 25.1 (Paso 1) ---
\begin{frame}{\small Ejemplo 1 - Paso 1}
\textbf{EDO:}
\[
  \frac{dy}{dx} = -2x^3 + 12x^2 - 20x + 8.5
\]
\textbf{Condiciones:}
\begin{itemize}
  \item $x=0$ a $x=4$ con $h=0.5$
  \item $y(0)=1$
\end{itemize}
\textbf{Paso 1:}
\begin{align*}
  f(0,1) &= -2(0)^3 + 12(0)^2 - 20(0) + 8.5 = 8.5 \\
  y(0.5) &= 1.0 + 8.5(0.5) = 5.25
\end{align*}
\textbf{Valor exacto:} $y(0.5) = 3.21875$ \hfill \textbf{Error:} $-2.03125$
\end{frame}

%--- Diapositiva: Ejemplo 25.1 (Paso 2) ---
\begin{frame}{\small  - Paso 2}
\textbf{Paso 2:}
\begin{align*}
  f(0.5, 5.25) &= -2(0.5)^3 + 12(0.5)^2 - 20(0.5) + 8.5 \\
  &= -0.25 + 3 -10 + 8.5 = 1.25 \\
  y(1) &= 5.25 + 0.5(1.25) = 5.875
\end{align*}
\textbf{Valor exacto:} $y(1) = 3.000$ \hfill \textbf{Error:} $-2.875$
\end{frame}

%--- Diapositiva: Tabla resumen del ejemplo ---
\begin{frame}{\small Resumen del Ejemplo 1}
\begin{center}
\begin{tabular}{|c|c|c|c|}
\hline
$x$ & $y_{\text{exacto}}$ & $y_{\text{Euler}}$ & Error \% \\
\hline
0.0 & 1.00000 & 1.00000 & 0.0 \\
0.5 & 3.21875 & 5.25000 & -63.1 \\
1.0 & 3.00000 & 5.87500 & -95.8 \\
1.5 & 2.21875 & 5.12500 & 131.0 \\
2.0 & 2.00000 & 4.50000 & -125.0 \\
2.5 & 2.71875 & 4.75000 & -74.7 \\
3.0 & 4.00000 & 5.87500 & 46.9 \\
3.5 & 4.71875 & 7.12500 & -51.0 \\
4.0 & 3.00000 & 7.00000 & -133.3 \\
\hline
\end{tabular}
\end{center}
\end{frame}

%--- Diapositiva: An\'alisis del error del m\'etodo de Euler ---
\begin{frame}{\small An\'alisis del Error del M\'etodo de Euler}
\begin{itemize}
  \item Existen dos tipos de errores en la soluci\'on num\'erica de EDO:
  \begin{itemize}
    \item \textbf{Errores de truncamiento (discretizaci\'on)}.
    \item \textbf{Errores de redondeo (precisi\'on finita)}.
  \end{itemize}
  \item El \textbf{error de truncamiento local} es el error cometido en un solo paso.
  \item El \textbf{error global} es la suma del error local y los errores propagados de pasos anteriores.
\end{itemize}
\end{frame}

%--- Diapositiva: Desarrollo con Serie de Taylor ---
\begin{frame}{\small Desarrollo con Serie de Taylor}
\begin{itemize}
  \item Sea \( y' = f(x, y) \)
  \item Expansi\'on de Taylor:
  \[
    y_{i+1} = y_i + y'_i h + \frac{y''_i}{2!}h^2 + \dots + \frac{y^{(n)}_i}{n!}h^n + R_n
  \]
  \item Sustituyendo derivadas sucesivas de $f$:
  \[
    y_{i+1} = y_i + f(x_i, y_i)h + \frac{f'(x_i, y_i)}{2!}h^2 + \dots + O(h^{n+1})
  \]
\end{itemize}
\end{frame}

%--- Diapositiva: Error de truncamiento local ---
\begin{frame}{\small Error de Truncamiento Local}
\begin{itemize}
  \item El m\'etodo de Euler utiliza:
  \[
    y_{i+1} = y_i + f(x_i, y_i)h
  \]
  \item Comparado con la serie de Taylor:
  \[
    E_t = \frac{f'(x_i, y_i)}{2!}h^2 + O(h^3)
  \]
  \item Entonces, el error local es:
  \[
    E_a = O(h^2)
  \]
  \item \textbf{Conclusi\'on:} Reducir el paso $h$ reduce r\'apidamente el error local.
\end{itemize}
\end{frame}

%--- Diapositiva: Ejemplo 25.2 - Estimaci\'on del Error ---
\begin{frame}{\small Ejemplo 25.2 - Estimaci\'on del Error de Euler}
\begin{itemize}
  \item EDO: \( \frac{dy}{dx} = -2x^3 + 12x^2 - 20x + 8.5 \)
  \item Derivadas necesarias:
  \begin{align*}
    f'(x) &= -6x^2 + 24x - 20 \\
    f''(x) &= -12x + 24 \\
    f^{(3)}(x) &= -12
  \end{align*}
  \item Evaluamos en $x=0$ con $h=0.5$:
  \begin{align*}
    E_{t,2} &= \frac{f'(0)}{2!} h^2 = \frac{-20}{2} (0.5)^2 = -2.5 \\
    E_{t,3} &= \frac{f''(0)}{3!} h^3 = \frac{24}{6} (0.5)^3 = 0.5 \\
    E_{t,4} &= \frac{f^{(3)}(0)}{4!} h^4 = \frac{-12}{24} (0.5)^4 = -0.03125
  \end{align*}
\end{itemize}
\end{frame}

%--- Diapositiva: Resultado del Error Total ---
\begin{frame}{\small Resultado del Error Total en el Paso 1}
\begin{itemize}
  \item Suma de errores de truncamiento:
  \[
    E_t = -2.5 + 0.5 - 0.03125 = -2.03125
  \]
  \item Coincide exactamente con el error obtenido en el Ejemplo 25.1 para $x=0.5$.
  \item \textbf{Conclusi\'on:} El error de orden 2 es dominante: $|E_{t,2}| > |E_{t,3}| > |E_{t,4}|$.
\end{itemize}
\end{frame}

%--- Diapositiva: Limitaciones del An\'alisis de Error con Taylor ---
\begin{frame}{\small Limitaciones del An\'alisis de Error con Taylor}
\begin{itemize}
  \item La serie de Taylor solo estima el \textbf{error de truncamiento local}.
  \item No proporciona el \textbf{error global} (acumulado por varios pasos).
  \item En problemas reales:
  \begin{itemize}
    \item Las soluciones exactas no se conocen.
    \item Las derivadas requeridas pueden ser complicadas o no anal\'iticas.
  \end{itemize}
  \item Aun con estas limitaciones, Taylor ayuda a entender el comportamiento de los m\'etodos.
\end{itemize}
\end{frame}

%--- Diapositiva: Conclusiones sobre el Error del M\'etodo de Euler ---
\begin{frame}{\small Conclusiones sobre el Error del M\'etodo de Euler}
\begin{itemize}
  \item El \textbf{error local} es proporcional a $h^2$, el \textbf{error global} es proporcional a $h$.
  \item Disminuir el tama\~no de paso $h$ reduce ambos errores.
  \item El m\'etodo de Euler es exacto si la soluci\'on es una funci\'on \textbf{lineal}.
  \item En general, el m\'etodo de Euler es un \textbf{m\'etodo de primer orden}.
  \item Los m\'etodos de orden $n$ son exactos para polinomios de grado $n$.
\end{itemize}
\end{frame}

%--- Diapositiva: An\'alisis del error del m\'etodo de Euler ---
\begin{frame}{\small An\'alisis del Error del M\'etodo de Euler}
\begin{itemize}
  \item Existen dos tipos de errores en la soluci\'on num\'erica de EDO:
  \begin{itemize}
    \item \textbf{Errores de truncamiento (discretizaci\'on)}.
    \item \textbf{Errores de redondeo (precisi\'on finita)}.
  \end{itemize}
  \item El \textbf{error de truncamiento local} es el error cometido en un solo paso.
  \item El \textbf{error global} es la suma del error local y los errores propagados de pasos anteriores.
\end{itemize}
\end{frame}

%--- Diapositiva: Desarrollo con Serie de Taylor ---
\begin{frame}{\small Desarrollo con Serie de Taylor}
\begin{itemize}
  \item Sea \( y' = f(x, y) \)
  \item Expansi\'on de Taylor:
  \[
    y_{i+1} = y_i + y'_i h + \frac{y''_i}{2!}h^2 + \dots + \frac{y^{(n)}_i}{n!}h^n + R_n
  \]
  \item Sustituyendo derivadas sucesivas de $f$:
  \[
    y_{i+1} = y_i + f(x_i, y_i)h + \frac{f'(x_i, y_i)}{2!}h^2 + \dots + O(h^{n+1})
  \]
\end{itemize}
\end{frame}

%--- Diapositiva: Error de truncamiento local ---
\begin{frame}{\small Error de Truncamiento Local}
\begin{itemize}
  \item El m\'etodo de Euler utiliza:
  \[
    y_{i+1} = y_i + f(x_i, y_i)h
  \]
  \item Comparado con la serie de Taylor:
  \[
    E_t = \frac{f'(x_i, y_i)}{2!}h^2 + O(h^3)
  \]
  \item Entonces, el error local es:
  \[
    E_a = O(h^2)
  \]
  \item \textbf{Conclusi\'on:} Reducir el paso $h$ reduce r\'apidamente el error local.
\end{itemize}
\end{frame}

%--- Diapositiva: Ejemplo 25.2 - Estimaci\'on del Error ---
\begin{frame}{\small Ejemplo 25.2 - Estimaci\'on del Error de Euler}
\begin{itemize}
  \item EDO: \( \frac{dy}{dx} = -2x^3 + 12x^2 - 20x + 8.5 \)
  \item Derivadas necesarias:
  \begin{align*}
    f'(x) &= -6x^2 + 24x - 20 \\
    f''(x) &= -12x + 24 \\
    f^{(3)}(x) &= -12
  \end{align*}
  \item Evaluamos en $x=0$ con $h=0.5$:
  \begin{align*}
    E_{t,2} &= \frac{f'(0)}{2!} h^2 = \frac{-20}{2} (0.5)^2 = -2.5 \\
    E_{t,3} &= \frac{f''(0)}{3!} h^3 = \frac{24}{6} (0.5)^3 = 0.5 \\
    E_{t,4} &= \frac{f^{(3)}(0)}{4!} h^4 = \frac{-12}{24} (0.5)^4 = -0.03125
  \end{align*}
\end{itemize}
\end{frame}

%--- Diapositiva: Resultado del Error Total ---
\begin{frame}{\small Resultado del Error Total en el Paso 1}
\begin{itemize}
  \item Suma de errores de truncamiento:
  \[
    E_t = -2.5 + 0.5 - 0.03125 = -2.03125
  \]
  \item Coincide exactamente con el error obtenido en el Ejemplo 25.1 para $x=0.5$.
  \item \textbf{Conclusi\'on:} El error de orden 2 es dominante: $|E_{t,2}| > |E_{t,3}| > |E_{t,4}|$.
\end{itemize}
\end{frame}

%--- Diapositiva: Limitaciones del An\'alisis de Error con Taylor ---
\begin{frame}{\small Limitaciones del An\'alisis de Error con Taylor}
\begin{itemize}
  \item La serie de Taylor solo estima el \textbf{error de truncamiento local}.
  \item No proporciona el \textbf{error global} (acumulado por varios pasos).
  \item En problemas reales:
  \begin{itemize}
    \item Las soluciones exactas no se conocen.
    \item Las derivadas requeridas pueden ser complicadas o no anal\'iticas.
  \end{itemize}
  \item Aun con estas limitaciones, Taylor ayuda a entender el comportamiento de los m\'etodos.
\end{itemize}
\end{frame}

%--- Diapositiva: Conclusiones sobre el Error del M\'etodo de Euler ---
\begin{frame}{\small Conclusiones sobre el Error del M\'etodo de Euler}
\begin{itemize}
  \item El \textbf{error local} es proporcional a $h^2$, el \textbf{error global} es proporcional a $h$.
  \item Disminuir el tama\~no de paso $h$ reduce ambos errores.
  \item El m\'etodo de Euler es exacto si la soluci\'on es una funci\'on \textbf{lineal}.
  \item En general, el m\'etodo de Euler es un \textbf{m\'etodo de primer orden}.
  \item Los m\'etodos de orden $n$ son exactos para polinomios de grado $n$.
\end{itemize}
\end{frame}

%--- Diapositiva: Serie de Taylor de orden superior ---
\begin{frame}{\small Serie de Taylor de Orden Superior}
\begin{itemize}
  \item Se puede mejorar Euler incluyendo t\'erminos de mayor orden:
  \[
    y_{i+1} = y_i + f(x_i, y_i) h + \frac{f'(x_i, y_i)}{2!} h^2
  \]
  \item El error de truncamiento aproximado ser\'ia:
  \[
    E_a = \frac{f''(x_i, y_i)}{3!} h^3
  \]
  \item \textbf{Ventaja:} mejora la precisi\'on.
  \item \textbf{Desventaja:} requiere derivadas m\'as complejas.
\end{itemize}
\end{frame}

%--- Diapositiva: Complejidad en derivadas de orden superior ---
\begin{frame}{\small Complejidad en Derivadas Superiores}
\begin{itemize}
  \item Para EDO dependientes de $x$ e $y$, se requiere la regla de la cadena:
  \[
    f' = \frac{\partial f}{\partial x} + \frac{\partial f}{\partial y} \cdot \frac{dy}{dx}
  \]
  \item La segunda derivada se vuelve m\'as compleja:
  \[
    f'' = \frac{\partial^2 f}{\partial x^2} + 2 \frac{\partial^2 f}{\partial x \partial y} \cdot \frac{dy}{dx} + \frac{\partial^2 f}{\partial y^2} \cdot \left( \frac{dy}{dx} \right)^2 + \frac{\partial f}{\partial y} \cdot \frac{d^2y}{dx^2}
  \]
  \item Esta complejidad limita la aplicaci\'on de Taylor de orden superior.
\end{itemize}
\end{frame}

%--- Diapositiva: Mejoras del m\'etodo de Euler ---
\begin{frame}{\small Mejoras del M\'etodo de Euler}
\begin{itemize}
  \item El error clave en el m\'etodo de Euler proviene de asumir que la pendiente inicial se mantiene constante.
  \item Se pueden mejorar los resultados usando m\'etodos que consideren la pendiente al final del intervalo.
  \item Uno de ellos es el \textbf{m\'etodo de Heun}, tambi\'en conocido como \textbf{Euler mejorado}.
  \item Estos m\'etodos pertenecen a la clase de \textbf{m\'etodos de Runge-Kutta}.
\end{itemize}
\end{frame}

%--- Diapositiva: M\'etodo de Heun (predictor) ---
\begin{frame}{\small M\'etodo de Heun: Paso Predictor}
\begin{itemize}
  \item Calcular la pendiente inicial:
  \[
    f(x_i, y_i)
  \]
  \item Predecir un valor intermedio:
  \[
    y^0_{i+1} = y_i + h f(x_i, y_i)
  \]
  \item Este paso es equivalente al m\'etodo de Euler y se conoce como \textbf{predictor}.
\end{itemize}
\end{frame}

%--- Diapositiva: M\'etodo de Heun (corrector) ---
\begin{frame}{\small M\'etodo de Heun: Paso Corrector}
\begin{itemize}
  \item Calcular la pendiente final con la predicci\'on:
  \[
    f(x_{i+1}, y^0_{i+1})
  \]
  \item Promediar las pendientes:
  \[
    f_{\text{prom}} = \frac{f(x_i, y_i) + f(x_{i+1}, y^0_{i+1})}{2}
  \]
  \item Calcular el valor corregido:
  \[
    y_{i+1} = y_i + h \cdot f_{\text{prom}}
  \]
\end{itemize}
\end{frame}

%--- Diapositiva: M\'etodo de Heun: Interpretaci\'on ---
\begin{frame}{\small M\'etodo de Heun: Interpretaci\'on}
\begin{itemize}
  \item El m\'etodo de Heun es un \textbf{m\'etodo predictor-corrector de un solo paso}.
  \item Se puede repetir la correcci\'on iterativamente:
  \[
    \varepsilon_a = \frac{y^{(j)}_{i+1} - y^{(j-1)}_{i+1}}{y^{(j)}_{i+1}} \cdot 100\%
  \]
  \item Iterar hasta que el error relativo aproximado \( \varepsilon_a \) sea menor a una tolerancia predefinida.
\end{itemize}
\end{frame}

%--- Diapositiva: Ejemplo 25.5 - Paso 1 ---
\begin{frame}{\small Ejemplo 25.5 - Heun, Paso 1}
\begin{itemize}
  \item EDO: \( \frac{dy}{dx} = 4e^{0.8x} - 0.5y \), con \( y(0) = 2 \), paso $h=1$.
  \item Calcular la pendiente inicial: \( f(0, 2) = 4 - 1 = 3 \)
  \item Paso predictor:
  \[
    y^0_1 = 2 + 1 \cdot 3 = 5
  \]
  \item Calcular \( f(1, 5) = 4e^{0.8} - 2.5 = 6.402164 \)
  \item Pendiente promedio:
  \[
    f_{prom} = \frac{3 + 6.402164}{2} = 4.701082
  \]
  \item Paso corrector:
  \[
    y_1 = 2 + 1 \cdot 4.701082 = 6.701082
  \]
\end{itemize}
\end{frame}

%--- Diapositiva: Ejemplo 25.5 - Iteraciones del corrector ---
\begin{frame}{\small Ejemplo 25.5 - Iteraciones del Corrector}
\begin{itemize}
  \item Segunda iteraci\'on:
  \[
    f(1, 6.701082) = 3.275811 \Rightarrow y_1^{(2)} = 2 + 1 \cdot \frac{3 + 3.275811}{2} = 6.275811
  \]
  \item Tercera iteraci\'on:
  \[
    f(1, 6.275811) = 3.382129 \Rightarrow y_1^{(3)} = 2 + \frac{3 + 3.382129}{2} = 6.382129
  \]
  \item Iterar hasta que:
  \[
    \varepsilon_a = \left| \frac{y_1^{(j)} - y_1^{(j-1)}}{y_1^{(j)}} \right| \cdot 100\% < \text{tolerancia}
  \]
  \item Con 15 iteraciones: \( y_1 \approx 6.360865 \) con error \( \approx 2.68\% \).
\end{itemize}
\end{frame}
%--- Diapositiva: Método de Heun para f(x) ---
\begin{frame}{\small Método de Heun cuando \( f = f(x) \)}
\begin{itemize}
  \item Si la EDO es \( \frac{dy}{dx} = f(x) \), no depende de \( y \).
  \item No se requiere predictor iterativo.
  \item Se aplica directamente:
  \[
    y_{i+1} = y_i + \frac{f(x_i) + f(x_{i+1})}{2} h
  \]
  \item Esta es la \textbf{regla del trapecio} aplicada a la integral de la EDO:
  \[
    y(x_{i+1}) = y(x_i) + \int_{x_i}^{x_{i+1}} f(x)dx
  \]
\end{itemize}
\end{frame}

%--- Diapositiva: Relación con la regla del trapecio ---
\begin{frame}{\small Relación con la Regla del Trapecio}
\begin{itemize}
  \item Partimos de:
  \[
    \frac{dy}{dx} = f(x) \Rightarrow y_{i+1} = y_i + \int_{x_i}^{x_{i+1}} f(x) dx
  \]
  \item Aplicando la regla del trapecio:
  \[
    \int_{x_i}^{x_{i+1}} f(x) dx \approx \frac{f(x_i) + f(x_{i+1})}{2} h
  \]
  \item Resulta en:
  \[
    y_{i+1} = y_i + \frac{f(x_i) + f(x_{i+1})}{2} h
  \]
  \item El error de truncamiento local:
  \[
    E_t = -\frac{f''(\xi)}{12} h^3, \quad \xi \in (x_i, x_{i+1})
  \]
\end{itemize}
\end{frame}
%--- Diapositiva: Métodos de Runge-Kutta ---
\begin{frame}{\small Métodos de Runge-Kutta}
\begin{itemize}
  \item Logran exactitud tipo Taylor sin derivadas de orden superior.
  \item Todos tienen la forma general:
  \[
    y_{i+1} = y_i + f(x_i, y_i, h) \cdot h
  \]
  \item Donde la \textbf{función incremento} es una combinación ponderada:
  \[
    f = a_1 k_1 + a_2 k_2 + \dots + a_n k_n
  \]
\end{itemize}
\end{frame}

%--- Diapositiva: Cálculo de los coeficientes k ---
\begin{frame}{\small Cálculo de los coeficientes \( k \) en Runge-Kutta}
\begin{itemize}
  \item Cada \( k \) representa una evaluación funcional:
  \begin{align*}
    k_1 &= f(x_i, y_i) \\
    k_2 &= f(x_i + p_1 h, y_i + q_{11} k_1 h) \\
    k_3 &= f(x_i + p_2 h, y_i + q_{21} k_1 h + q_{22} k_2 h) \\
    \vdots \\
    k_n &= f(x_i + p_{n-1} h, y_i + \sum q_{n-1,j} k_j h)
  \end{align*}
  \item Los valores \( a, p, q \) se eligen para coincidir con la serie de Taylor.
\end{itemize}
\end{frame}

%--- Diapositiva: Orden y precisión ---
\begin{frame}{\small Orden y Precisión de los Métodos RK}
\begin{itemize}
  \item RK de orden 1: equivalente al método de Euler.
  \item RK de orden 2: como el método del punto medio o Heun.
  \item RK de orden 3 y 4: mayor precisión sin derivadas de orden superior.
  \item Errores:
  \begin{itemize}
    \item Local: \( O(h^{n+1}) \)
    \item Global: \( O(h^n) \)
  \end{itemize}
\end{itemize}
\end{frame}
%--- Diapositiva: RK de segundo orden ---
\begin{frame}{\small Runge-Kutta de Segundo Orden}
\begin{itemize}
  \item Fórmula general:
  \[
    y_{i+1} = y_i + (a_1 k_1 + a_2 k_2) h
  \]
  \item Donde:
  \[
    k_1 = f(x_i, y_i), \quad k_2 = f(x_i + p_1 h, y_i + q_{11} k_1 h)
  \]
  \item Se elige un valor para una constante (por ejemplo, \( a_2 \)) y se despejan las otras para hacer coincidir con la expansión de Taylor.
\end{itemize}
\end{frame}

%--- Diapositiva: Versiones comunes de RK2 ---
\begin{frame}{\small Versiones comunes de RK de Segundo Orden}
\begin{itemize}
  \item \textbf{Heun (}\( a_2 = 1/2 \)): \( a_1 = 1/2,\; p_1 = q_{11} = 1 \)
  \[
    y_{i+1} = y_i + \frac{1}{2}(k_1 + k_2) h
  \]
  \item \textbf{Punto medio (}\( a_2 = 1 \)): \( a_1 = 0,\; p_1 = q_{11} = 1/2 \)
  \[
    y_{i+1} = y_i + k_2 h
  \]
  \item \textbf{Ralston (}\( a_2 = 2/3 \)): \( a_1 = 1/3,\; p_1 = q_{11} = 3/4 \)
  \[
    y_{i+1} = y_i + \left(\frac{1}{3}k_1 + \frac{2}{3}k_2\right) h
  \]
\end{itemize}
\end{frame}

%--- Diapositiva: Ejemplo 25.6 Punto Medio y Ralston ---
\begin{frame}{\small Ejemplo 25.6: Comparación de Métodos RK2}
\textbf{EDO:} \( f(x) = -2x^3 + 12x^2 - 20x + 8.5 \)
\\
\textbf{Condición inicial:} \( y(0) = 1 \), \( h = 0.5 \), integrar hasta \( x = 4 \)
\begin{itemize}
  \item Punto medio:
  \[ k_1 = f(0) = 8.5 \quad \Rightarrow \quad k_2 = f(0.25) = 4.21875 \]
  \[ y_1 = 1 + 4.21875 \cdot 0.5 = 3.109375 \]
  \item Ralston:
  \[ k_1 = 8.5 \quad \Rightarrow \quad k_2 = f(0.375) = 2.58203125 \]
  \[ y_1 = 1 + \left(\frac{1}{3}\cdot 8.5 + \frac{2}{3}\cdot 2.58203125\right) \cdot 0.5 \approx 3.2773 \]
\end{itemize}
\end{frame}

%--- Diapositiva: Resultados Comparativos ---
\begin{frame}{\small Comparación de Resultados: Tabla 25.3}
\scriptsize
\begin{tabular}{|c|c|c|c|c|c|c|c|}
\hline
$x$ & $y_{\text{verdadero}}$ & Heun & $\%\varepsilon_t$ & Punto medio & $\%\varepsilon_t$ & Ralston & $\%\varepsilon_t$ \\
\hline
0.0 & 1.0000 & 1.0000 & 0 & 1.0000 & 0 & 1.0000 & 0 \\
0.5 & 3.2188 & 3.4375 & 6.8 & 3.1094 & 3.4 & 3.2773 & 1.8 \\
1.0 & 3.0000 & 3.3750 & 12.5 & 2.8125 & 6.3 & 3.1016 & 3.4 \\
1.5 & 2.2188 & 2.6875 & 21.1 & 1.9844 & 10.6 & 2.3477 & 5.8 \\
2.0 & 2.0000 & 2.5000 & 25.0 & 1.7500 & 12.5 & 2.1406 & 7.0 \\
2.5 & 2.7188 & 3.1875 & 17.2 & 2.4844 & 8.6 & 2.8555 & 5.0 \\
3.0 & 4.0000 & 4.3750 & 9.4 & 3.8125 & 4.7 & 4.1172 & 2.9 \\
3.5 & 4.7188 & 4.9375 & 4.6 & 4.6094 & 2.3 & 4.8008 & 1.7 \\
4.0 & 3.0000 & 3.0000 & 0 & 3.0000 & 0 & 3.0313 & 1.0 \\
\hline
\end{tabular}
\end{frame}

%--- Diapositiva: RK de tercer orden ---
\begin{frame}{\small Runge-Kutta de Tercer Orden}
\begin{itemize}
  \item Forma común:
  \[
    y_{i+1} = y_i + \frac{1}{6}(k_1 + 4k_2 + k_3) h
  \]
  \item Donde:
  \begin{align*}
    k_1 &= f(x_i, y_i) \\
    k_2 &= f\left(x_i + \frac{1}{2}h, y_i + \frac{1}{2}k_1 h\right) \\
    k_3 &= f(x_i + h, y_i - k_1 h + 2 k_2 h)
  \end{align*}
  \item Exacto para ecuaciones cúbicas cuando la solución es polinómica de orden ≤ 3.
  \item Si \( f = f(x) \), equivale a la regla de Simpson 1/3.
\end{itemize}
\end{frame}

%--- Diapositiva: RK de cuarto orden ---
\begin{frame}{\small Runge-Kutta de Cuarto Orden (Clásico)}
\begin{itemize}
  \item Fórmula:
  \[
    y_{i+1} = y_i + \frac{1}{6}(k_1 + 2k_2 + 2k_3 + k_4) h
  \]
  \item Donde:
  \begin{align*}
    k_1 &= f(x_i, y_i) \\
    k_2 &= f\left(x_i + \frac{h}{2}, y_i + \frac{h}{2}k_1\right) \\
    k_3 &= f\left(x_i + \frac{h}{2}, y_i + \frac{h}{2}k_2\right) \\
    k_4 &= f(x_i + h, y_i + h k_3)
  \end{align*}
  \item Promedio ponderado de pendientes.
  \item Similitud con Heun, pero con mayor precisión: \( O(h^5) \) local, \( O(h^4) \) global.
\end{itemize}
\end{frame}
%--- Diapositiva: Ejemplo RK4 clásico ---
\begin{frame}{\small Ejemplo RK4 clásico: Caso (a)}
\textbf{EDO:} \( f(x, y) = -2x^3 + 12x^2 - 20x + 8.5 \),\quad \( h = 0.5,\; y(0) = 1 \)
\begin{itemize}
  \item \( k_1 = f(0,1) = 8.5 \)
  \item \( k_2 = f(0.25, 2.125) = 4.21875 \)
  \item \( k_3 = f(0.25, 2.0546875) = 4.21875 \)
  \item \( k_4 = f(0.5, 3.0546875) = 1.25 \)
  \item Resultado:
  \[
    y(0.5) = 1 + \frac{1}{6}(8.5 + 2\cdot4.21875 + 2\cdot4.21875 + 1.25)(0.5) = 3.21875
  \]
\end{itemize}
\end{frame}

%--- Diapositiva: Ejemplo RK4 clásico: Caso (b) ---
\begin{frame}{\small Ejemplo RK4 clásico: Caso (b)}
\textbf{EDO:} \( f(x, y) = 4e^{0.8x} - 0.5y \),\quad \( h = 0.5,\; y(0) = 2 \)
\begin{itemize}
  \item \( k_1 = 4e^{0} - 0.5(2) = 3 \)
  \item \( k_2 = 4e^{0.2} - 0.5(2.75) = 3.510611 \)
  \item \( k_3 = 4e^{0.2} - 0.5(2.877653) = 3.446785 \)
  \item \( k_4 = 4e^{0.5} - 0.5(3.723392) = 4.105603 \)
  \item Resultado:
  \[
    y(0.5) = 2 + \frac{1}{6}(3 + 2\cdot3.510611 + 2\cdot3.446785 + 4.105603)(0.5) = 3.751669
  \]
\end{itemize}
\end{frame}
%--- Diapositiva: RK de cuarto orden ---
\begin{frame}{\small Runge-Kutta de Cuarto Orden (Clásico)}
\begin{itemize}
  \item Fórmula:
  \[
    y_{i+1} = y_i + \frac{1}{6}(k_1 + 2k_2 + 2k_3 + k_4) h
  \]
  \item Donde:
  \begin{align*}
    k_1 &= f(x_i, y_i) \\
    k_2 &= f\left(x_i + \frac{h}{2}, y_i + \frac{h}{2}k_1\right) \\
    k_3 &= f\left(x_i + \frac{h}{2}, y_i + \frac{h}{2}k_2\right) \\
    k_4 &= f(x_i + h, y_i + h k_3)
  \end{align*}
  \item Promedio ponderado de pendientes.
  \item Similitud con Heun, pero con mayor precisión: \( O(h^5) \) local, \( O(h^4) \) global.
\end{itemize}
\end{frame}

%--- Diapositiva: RK de quinto orden de Butcher ---
\begin{frame}{\small Runge-Kutta de Quinto Orden (Butcher)}
\begin{itemize}
  \item Fórmula:
  \[
    y_{i+1} = y_i + \frac{1}{90}(7k_1 + 32k_3 + 12k_4 + 32k_5 + 7k_6) h
  \]
  \item Coeficientes:
  \begin{align*}
    k_1 &= f(x_i, y_i) \\
    k_2 &= f\left(x_i + \frac{1}{4}h, y_i + \frac{1}{4}k_1 h\right) \\
    k_3 &= f\left(x_i + \frac{1}{4}h, y_i + \frac{1}{8}k_1 h + \frac{1}{8}k_2 h\right) \\
    k_4 &= f\left(x_i + \frac{1}{2}h, y_i - \frac{1}{2}k_2 h + k_3 h\right) \\
    k_5 &= f\left(x_i + \frac{3}{4}h, y_i + \frac{3}{16}k_1 h + \frac{9}{16}k_4 h\right) \\
    k_6 &= f\left(x_i + h, y_i - \frac{3}{7}k_1 h + \frac{2}{7}k_2 h + \frac{12}{7}k_3 h - \frac{12}{7}k_4 h + \frac{8}{7}k_5 h\right)
  \end{align*}
  \item Mayor precisión pero con más evaluaciones de la función.
\end{itemize}
\end{frame}
\end{document}